\documentclass[9pt]{beamer-control}
\usepackage{beamer-control-prac}
\begin{document}
\CONCEPT[7]{Practical Project: Ping Pong control system}

\begin{frame}
\frametitle{Introduction}
In this practical project, you will apply control theory concepts to design and implement a control system for a complex ping pong ball manipulation rig.

\vfill

This practical project will consist of the following parts:
\begin{itemize}
\item System identification and modelling
\item Controller design and tuning
\item Performance evaluation and analysis
\end{itemize}
You will be assessed via a formal group report.

\end{frame}

\begin{frame}
\frametitle{Timing}

\begin{itemize}
\item Each part should take approximately \emph{no more than} one week to complete.
\item Plan to write up the report as you go, as it will also take time to put together.
\item Manage your time and progress carefully. A final (fourth) prac session will be held in the final week of semester to allow for finishing up, but no sessions are planned for SWOTVAC.
\end{itemize}
\end{frame}

\SUBCONCEPT{Problem statement}

\begin{frame}{Problem statement}
In this project, you are provided with a hardware rig that needs to be controlled.
\begin{itemize}
\item Using your best engineering judgement, you are required to implement a position controller for the height of the ping pong ball
\item As part of the documentation, a complete description of the hardware and interfaces is required
\item In order to achieve setpoint following, calibration of the position sensor is required
\item No fixed performance criteria is specified; in your documentation you must justify how you tuned the controller and attempted to optimise the performance of the closed loop system
\item While a relatively simple PID controller in its `pure' form will manage to control this plant, we encourage creativity in terms of signal processing, fudge factors, gain scheduling, in order to improve performance
\end{itemize}
\end{frame}

\SUBCONCEPT{Part A: System identification}

\begin{frame}{Interfacing}
\begin{itemize}
\item Identifying inputs and outputs
\item Calibration of sensor input signals
\item Bounding output signals
\end{itemize}
\end{frame}


\begin{frame}{Modelling the ping pong rig}
\begin{itemize}
\item Identify system dynamics
\item Determine transfer function and/or state space model
\end{itemize}
\end{frame}


\SUBCONCEPT{Part B: Controller design and tuning}

\begin{frame}{Implementing control strategies}
\begin{itemize}
\item Design appropriate controller(s) --- more req'd than trial-and-error
\item Measure initial performance characteristics
\end{itemize}
\end{frame}


\SUBCONCEPT{Part C: Performance evaluation and analysis}

\begin{frame}{Performance evaluation}
\begin{itemize}
\item Use systematic methods to improve the tuning
\item Quantify performance improvements against at least two metrics (e.g., overshoot, \dots)
\end{itemize}
\end{frame}


\begin{frame}{Analysis and testing}
\begin{itemize}
\item Test against a sweep of step inputs (different amplitudes, different signs)
\item Use at least one approach to test disturbance rejection characteristics
\end{itemize}
\end{frame}


\SUBCONCEPT{Report and assessment}

\begin{frame}{Rubric}
The project is assessed via a formal report. The rubric includes the following criteria (equally weighted):

\small
\begin{enumerate}
    \item \textbf{Writing Style and Report Formatting}:
    Grammar, spelling, tenses, formality, completeness, flow. Etc. Appropriate introduction and conclusion. Must include executive summary. High quality figures and graphs, typography, typesetting, document structure, tables of contents, figure captions. Etc.
    
    \item \textbf{Systems overview}:
    Explanation (and possibly derivation) of dynamics, actuator, interfacing with sensor(s), input signals, output signals, ranges of electrical signals. Photo(s) / schematic(s) / flow diagram(s) needed.

    \item \textbf{Open loop characterisations}:
    Sensor calibration, signal normalisation, step responses, nonlinearities, frequency responses. Characterisation of plant model (transfer function, ITD/FOTD parameters, etc., as appropriate).

    \item \textbf{Control systems}:
    Use of at least one tuning recipe, characterising performance (settling time, overshoot, etc.). Manual tuning improvements to improve default response. Discussion of desired performance and obstacles for improving performance.

    \item \textbf{Performance and robustness}:
    Exploration of system behaviour under varying conditions, including but not limited to robustness to parameter variation (e.g., mass variation), measurement noise, disturbances. Changes in performance vs variations. Limits of stability.
\end{enumerate}
\end{frame}

\SUMMARYFRAME

\end{document}
